% =====================================================================
%  Preambolo condiviso per i documenti d'esame (classe exam)
%  Sicurezza delle Applicazioni e IA per la Cybersecurity - Modulo 1
%  UniParthenope - DISEGIM - A.A. 2026/2027 - Dr. Andrea Bruno
%  babel italiano caricato solo se disponibile (IfFileExists).
% =====================================================================
\usepackage[T1]{fontenc}
\usepackage[utf8]{inputenc}
\IfFileExists{italian.ldf}{\usepackage[italian]{babel}}{}
\usepackage{lmodern}
\usepackage{amsmath,amssymb}
\usepackage{xcolor}
\usepackage{listings}
\usepackage{enumitem}
\usepackage{booktabs}
\usepackage{array}
\usepackage[hidelinks]{hyperref}
\definecolor{codebg}{HTML}{F6F8FA}
\definecolor{codekw}{HTML}{0033B3}
\definecolor{codecm}{HTML}{6A737D}
\definecolor{codestr}{HTML}{067D17}
\lstset{
  basicstyle=\small\ttfamily,
  backgroundcolor=\color{codebg},
  keywordstyle=\color{codekw}\bfseries,
  commentstyle=\color{codecm}\itshape,
  stringstyle=\color{codestr},
  numbers=left, numberstyle=\tiny\color{codecm}, numbersep=6pt,
  breaklines=true, showstringspaces=false,
  frame=single, rulecolor=\color{codecm!40},
  extendedchars=true, inputencoding=utf8,
  literate={à}{{\`a}}1 {è}{{\`e}}1 {é}{{\'e}}1 {ì}{{\`i}}1 {ò}{{\`o}}1 {ù}{{\`u}}1 {È}{{\`E}}1 {°}{{\textdegree}}1,
  aboveskip=0.6em, belowskip=0.6em
}
\lstdefinestyle{c}{language=C}
\lstdefinestyle{py}{language=Python}
\lstdefinestyle{sh}{language=bash, numbers=none}
\lstdefinestyle{asm}{language={[x86masm]Assembler}, numbers=none}
\providecommand{\titolovariante}{}
\pagestyle{headandfoot}
\runningheadrule
\firstpageheader{\small Sic. App.\ -- Mod.\,1}{}{\small A.A.\ 2026/2027}
\runningheader{\small Sic. App.\ -- Mod.1}{\small \lezionenome}{\small \thepage}
\runningfooter{}{}{}
\pointsinrightmargin
\marginpointname{\,pt}
\newcommand{\intestazione}{%
  \begin{center}
    {\Large\bfseries Sicurezza delle Applicazioni e IA per la Cybersecurity -- Modulo 1}\\[2pt]
    {\large \lezionenome}\\[3pt]
    {\normalsize Esercitazione \titolovariante}\\[6pt]
    {\footnotesize Universit\`a degli Studi di Napoli ``Parthenope'' -- DISEGIM -- A.A.\ 2026/2027 -- Dr.\ Andrea Bruno}
  \end{center}
  \vspace{2pt}\hrule\vspace{6pt}
  \noindent\makebox[\textwidth][s]{\strut Cognome e nome: \underline{\hspace{6cm}}\hfill Matricola: \underline{\hspace{3.5cm}}}\\[6pt]
  \noindent Totale punti: \numpoints.\quad Tempo: 90 minuti.\quad Motivare sempre le risposte.
  \vspace{4pt}\hrule\vspace{12pt}
}

% ---- spazio per la risposta ----
% Nell'esercitazione (senza soluzioni) il corpo di ogni ambiente
% solution viene scartato e al suo posto restano delle righe vuote su cui
% scrivere; con \printanswers (soluzioni.tex) compare la soluzione.
\usepackage{comment}
\newcommand{\altezzarisposta}{5cm}
\ifprintanswers\else
  \let\solution\relax\let\endsolution\relax
  \excludecomment{solution}
  \let\saltasoluzione\solution
  \def\solution{\par\fillwithlines{\altezzarisposta}\saltasoluzione}
\fi

% ---- \fullwidth prima della prima \question ----
% Le versioni recenti di TeX Live segnalano "missing \item" se \fullwidth
% apre un paragrafo prima del primo \item: azzera il flag di lista nuova.
\makeatletter
\let\oldfullwidth\fullwidth
\renewcommand{\fullwidth}[1]{\@newlistfalse\oldfullwidth{#1}}
\makeatother
